\chapter{2010年辽宁卷（理）}

\section{单选题}

\begin{problem}
已知 \(A\)，\(B\) 均为集合 \(U = \{1, 3, 5, 7, 9\}\) 的子集，且 \(A \cap B = \{3\}\)，\(\complement_U B \cap A = \{9\}\)，则 \(A =\)（\quad）

\choices
    {\(\{1, 3\}\)}
    {\(\{3,7,9\}\)}
    {\(\{3,5,9\}\)}
    {\(\{3,9\}\)}
\end{problem}

\begin{answer}
D
\end{answer}

\begin{solution}
由于 \(U\) 中的任意元素不是属于 \(B\)，就是属于 \(\complement_U B\)，所以
\[
A=(A\cap B)\cup(A\cap\complement_U B)=\{3\}\cup\{9\}=\{3,9\}
\]
故选 D.
\end{solution}

\begin{problem}
设 \(a\)，\(b\) 为实数，若复数 \(\frac{1 + 2\i}{a + \i b} = 1 + \i\)，则（\quad）

\choices
    {\(a = \frac{3}{2}, b = \frac{1}{2}\)}
    {\(a = 3\)，\(b = 1\)}
    {\(a = \frac{1}{2}\)，\(b = \frac{3}{2}\)}
    {\(a = 1\)，\(b = 3\)}
\end{problem}

\begin{answer}
A
\end{answer}

\begin{solution}
由题意，分母 \(a+\i b\neq0\)，交叉相乘得
\[
1+2\i=(a+\i b)(1+\i)=(a-b)+(a+b)\i
\]
比较实部和虚部，得到 \(a-b=1\)，\(a+b=2\). 两式相加、相减，分别得 \(a=\frac{3}{2}\)，\(b=\frac{1}{2}\)，故选 A.
\end{solution}

\begin{problem}
两个实习生每人加工一个零件，加工为一等品的概率分别为 \(\frac{2}{3}\) 和 \(\frac{3}{4}\)，两个零件是否加工为一等品相互独立，则这两个零件中恰有一个一等品的概率为（\quad）

\choices
    {\( \frac{1}{2} \)}
    {\( \frac{5}{12} \)}
    {\( \frac{1}{4} \)}
    {\( \frac{1}{6} \)}
\end{problem}

\begin{answer}
B
\end{answer}

\begin{solution}
设两名实习生加工的零件分别为 \(X\)，\(Y\)，则恰有一个一等品包含两个互斥事件：\(X\) 是一等品且 \(Y\) 不是一等品，或 \(X\) 不是一等品且 \(Y\) 是一等品. 由独立性，所求概率为
\[
\frac{2}{3}\left(1-\frac{3}{4}\right)+\left(1-\frac{2}{3}\right)\frac{3}{4}=\frac{2}{3}\cdot\frac{1}{4}+\frac{1}{3}\cdot\frac{3}{4}=\frac{5}{12}
\]
故选 B.
\end{solution}

\begin{problem}
如果执行如图所示的程序框图，输入正整数 \(n\)，\(m\)，满足 \(n \geqslant m\)，那么输出的 \(p\) 等于（\quad）

\texfigure[max width=0.34\linewidth]{img_repaint/2010/liaoning_science/q4_fig1.tex}

\choices
    {\(\mathrm C_n^{m-1}\)}
    {\(\mathrm A_n^{m-1}\)}
    {\(\mathrm C_n^m\)}
    {\(\mathrm A_n^m\)}
\end{problem}

\begin{answer}
D
\end{answer}

\begin{solution}
程序开始时 \(k=1\)，\(p=1\). 每次循环先令 \(p=p(n-m+k)\)，再令 \(k=k+1\)；当 \(k<m\) 时继续循环. 因此循环中 \(k\) 依次取 \(1,2,\ldots,m\)，最终
\[
p=\prod_{k=1}^{m}(n-m+k)=(n-m+1)(n-m+2)\cdots n=\frac{n!}{(n-m)!}=\mathrm A_n^m
\]
故选 D.
\end{solution}

\begin{problem}
设 \(\omega > 0\)，函数 \(y = \sin \left(\omega x + \frac{\pi}{3}\right) + 2\) 的图象向右平移 \(\frac{4\pi}{3}\) 个单位后与原图象重合，则 \(\omega\) 的最小值是（\quad）

\choices
    {\(\frac{2}{3}\)}
    {\( \frac{4}{3} \)}
    {\( \frac{3}{2} \)}
    {\(3\)}
\end{problem}

\begin{answer}
C
\end{answer}

\begin{solution}
设原函数为 \(f(x)=\sin\left(\omega x+\frac{\pi}{3}\right)+2\). 图象向右平移 \(\frac{4\pi}{3}\) 个单位后得到
\[
f\left(x-\frac{4\pi}{3}\right)=\sin\left(\omega x+\frac{\pi}{3}-\frac{4\pi\omega}{3}\right)+2
\]
它与原图象重合，需满足相位差 \(\frac{4\pi\omega}{3}\) 是 \(2\pi\) 的正整数倍，即 \(\frac{4\pi\omega}{3}=2k\pi\)，其中 \(k\in\mathbb Z_{+}\).
所以 \(\omega=\frac{3k}{2}\)，其最小值为 \(\frac{3}{2}\)，故选 C.
\end{solution}

\begin{problem}
设 \(\{a_{n}\}\) 是由正数组成的等比数列，\(S_{n}\) 为其前 \(n\) 项和. 已知 \(a_{2}a_{4} = 1\)，\(S_{3} = 7\)，则 \(S_{5} =\)（\quad）

\choices
    {\( \frac{15}{2} \)}
    {\( \frac{31}{4} \)}
    {\( \frac{33}{4} \)}
    {\( \frac{17}{2} \)}
\end{problem}

\begin{answer}
B
\end{answer}

\begin{solution}
设等比数列的公比为 \(r>0\). 由等比数列性质，\(a_{2}a_{4}=a_{3}^{2}\)，又因各项为正数，所以 \(a_{3}=1\). 于是
\(a_{2}=\frac{1}{r}\)，\(a_{1}=\frac{1}{r^{2}}\). 由 \(S_{3}=7\) 得
\[
\frac{1}{r^{2}}+\frac{1}{r}+1=7
\]
即 \(6r^{2}-r-1=0\).
解得 \(r=\frac{1}{2}\) 或 \(r=-\frac{1}{3}\)，结合 \(r>0\)，得 \(r=\frac{1}{2}\). 因此
\(a_{1}=4\)，\(a_{2}=2\)，\(a_{3}=1\)，\(a_{4}=\frac{1}{2}\)，\(a_{5}=\frac{1}{4}\).
从而 \(S_{5}=4+2+1+\frac{1}{2}+\frac{1}{4}=\frac{31}{4}\)，故选 B.
\end{solution}

\begin{problem}
设抛物线 \( y^{2} = 8x \) 的焦点为 \(F\)，准线为 \( l \)，\(P\) 为抛物线上一点，\( PA \perp l \)，\(A\) 为垂足，如果直线 \(AF\) 的斜率为 \( -\sqrt{3} \)，那么 \( |PF| = \)（\quad）

\choices
    {\(4\sqrt{3}\)}
    {\(8\)}
    {\( 8\sqrt{3} \)}
    {\(16\)}
\end{problem}

\begin{answer}
B
\end{answer}

\begin{solution}
将抛物线写成 \(y^{2}=4px\)，得 \(p=2\)，所以焦点为 \(F=(2,0)\)，准线为 \(l:x=-2\). 由于 \(PA\perp l\)，线段 \(PA\) 水平，设 \(P=(x,y)\)，则 \(A=(-2,y)\). 由直线 \(AF\) 的斜率为 \(-\sqrt{3}\)，得
\[
\frac{0-y}{2-(-2)}=-\sqrt{3}
\]
所以 \(y=4\sqrt{3}\).
又 \(P\) 在抛物线上，所以 \(x=\frac{y^{2}}{8}=6\). 因此
\[
|PF|=\sqrt{(6-2)^{2}+(4\sqrt{3})^{2}}=\sqrt{16+48}=8
\]
故选 B.
\end{solution}

\begin{problem}
平面上 \(O\)，\(A\)，\(B\) 三点不共线，设 \(\overrightarrow{OA} = \bs{a}\)，\(\overrightarrow{OB} = \bs{b}\)，则 \(\triangle OAB\) 的面积等于（\quad）

\choices
    {\(\sqrt{|\bs{a}|^2|\bs{b}|^2 - (\bs{a} \cdot \bs{b})^2}\)}
    {\(\sqrt{|\bs{a}|^2|\bs{b}|^2 + (\bs{a} \cdot \bs{b})^2}\)}
    {\(\frac{1}{2}\sqrt{|\bs{a}|^2|\bs{b}|^2 - (\bs{a} \cdot \bs{b})^2}\)}
    {\(\frac{1}{2}\sqrt{|\bs{a}|^2|\bs{b}|^2 + (\bs{a} \cdot \bs{b})^2}\)}
\end{problem}

\begin{answer}
C
\end{answer}

\begin{solution}
设向量 \(\bs{a}\)，\(\bs{b}\) 的夹角为 \(\theta\)，因为 \(O\)，\(A\)，\(B\) 不共线，所以 \(0<\theta<\pi\). 三角形面积为
\[
S_{\triangle OAB}=\frac{1}{2}|\bs{a}|\,|\bs{b}|\sin\theta
\]
另一方面，\(\bs{a}\cdot\bs{b}=|\bs{a}|\,|\bs{b}|\cos\theta\)，故
\[
\sqrt{|\bs{a}|^{2}|\bs{b}|^{2}-(\bs{a}\cdot\bs{b})^{2}}=|\bs{a}|\,|\bs{b}|\sin\theta
\]
所以三角形面积等于
\[
\frac{1}{2}\sqrt{|\bs{a}|^{2}|\bs{b}|^{2}-(\bs{a}\cdot\bs{b})^{2}}
\]
故选 C.
\end{solution}

\begin{problem}
设双曲线的一个焦点为 \(F\)，虚轴的一个端点为 \(B\)，如果直线 \(FB\) 与该双曲线的一条渐近线垂直，那么此双曲线的离心率为（\quad）

\choices
    {\(\sqrt{2}\)}
    {\(\sqrt{3}\)}
    {\(\frac{\sqrt{3} + 1}{2}\)}
    {\(\frac{\sqrt{5} + 1}{2}\)}
\end{problem}

\begin{answer}
D
\end{answer}

\begin{solution}
设双曲线方程为 \(\frac{x^{2}}{a^{2}}-\frac{y^{2}}{b^{2}}=1\)，其中 \(a>0\)，\(b>0\)，半焦距为 \(c\)，则 \(c^{2}=a^{2}+b^{2}\)，离心率为 \(e=\frac{c}{a}\). 取焦点 \(F=(c,0)\) 和虚轴端点 \(B=(0,b)\)，则直线 \(FB\) 的斜率为 \(-\frac{b}{c}\)，双曲线的一条渐近线斜率为 \(\frac{b}{a}\). 由两直线垂直，得
\[
\left(-\frac{b}{c}\right)\frac{b}{a}=-1
\]
所以 \(b^{2}=ac\).
两边除以 \(a^{2}\)，得 \(\frac{b^{2}}{a^{2}}=\frac{c}{a}=e\). 再由 \(c^{2}=a^{2}+b^{2}\)，得
\[
e^{2}=1+\frac{b^{2}}{a^{2}}=1+e
\]
即 \(e^{2}-e-1=0\).
因为 \(e>1\)，所以 \(e=\frac{1+\sqrt{5}}{2}\)，故选 D.
\end{solution}

\begin{problem}
已知点 \(P\) 在曲线 \(y = \frac{4}{\e^{x} + 1}\) 上，\(\alpha\) 为曲线在点 \(P\) 处的切线的倾斜角，则 \(\alpha\) 的取值范围是（\quad）

\choices
    {\(\left[0, \frac{\pi}{4}\right)\)}
    {\(\left[\frac{\pi}{4}, \frac{\pi}{2}\right)\)}
    {\(\left(\frac{\pi}{2}, \frac{3\pi}{4}\right]\)}
    {\(\left[\frac{3\pi}{4}, \pi\right)\)}
\end{problem}

\begin{answer}
D
\end{answer}

\begin{solution}
函数的导数为
\[
y'=-\frac{4\e^{x}}{(\e^{x}+1)^{2}}
\]
令 \(t=\e^{x}>0\)，则
\[
|y'|=\frac{4t}{(t+1)^{2}}\leqslant1
\]
其中不等式由 \((t-1)^{2}\geqslant0\) 得到，且 \(t=1\) 时取等号. 因此切线斜率满足 \(-1\leqslant y'<0\). 切线倾斜角 \(\alpha\) 在 \((\frac{\pi}{2},\pi)\) 内，且 \(\tan\alpha=y'\)，所以
\[
\frac{3\pi}{4}\leqslant\alpha<\pi
\]
故选 D.
\end{solution}

\begin{problem}
已知 \(a > 0\)，则 \(x_0\) 满足关于 \(x\) 的方程 \(ax = b\) 的充要条件是（\quad）

\choices
    {\(\exists x \in \R,\ \frac{1}{2}ax^{2} - bx \geqslant \frac{1}{2}ax_{0}^{2} - bx_{0}\)}
    {\(\exists x \in \R\)，\(\ \frac{1}{2}ax^{2} - bx \leqslant \frac{1}{2}ax_{0}^{2} - bx_{0}\)}
    {\(\forall x \in \R\)，\(\ \frac{1}{2}ax^{2} - bx \geqslant \frac{1}{2}ax_{0}^{2} - bx_{0}\)}
    {\(\forall x \in \R\)，\(\ \frac{1}{2}ax^{2} - bx \leqslant \frac{1}{2}ax_{0}^{2} - bx_{0}\)}
\end{problem}

\begin{answer}
C
\end{answer}

\begin{solution}
令
\[
g(x)=\frac{1}{2}ax^{2}-bx
\]
因为 \(a>0\)，配方得
\[
g(x)=\frac{a}{2}\left(x-\frac{b}{a}\right)^{2}-\frac{b^{2}}{2a}
\]
所以 \(g(x)\) 在 \(x=\frac{b}{a}\) 处取得唯一最小值. 于是
\[
ax_{0}=b
\]
当且仅当 \(\forall x\in\R\)，\(g(x)\geqslant g(x_{0})\)，这正是选项 C，故选 C.
\end{solution}

\begin{problem}
有四根长都为 2 的直铁条，若再选两根长都为 \(a\) 的直铁条，使这六根铁条端点处相连能够焊接成一个三棱锥形的铁架，则 \(a\) 的取值范围是（\quad）

\choices
    {\((0, \sqrt{6} + \sqrt{2})\)}
    {\((1, 2\sqrt{2})\)}
    {\((\sqrt{6} -\sqrt{2},\sqrt{6} +\sqrt{2})\)}
    {\((0,2\sqrt{2})\)}
\end{problem}

\begin{answer}
A
\end{answer}

\begin{solution}
六根铁条构成三棱锥的六条棱时，两根长为 \(a\) 的棱只有两种相对位置：相交于同一顶点，或为一对异棱.

先考虑两根长为 \(a\) 的棱相交于顶点 \(A\). 设 \(AB=AC=a\)，其余四条棱 \(AD\)，\(BC\)，\(BD\)，\(CD\) 长均为 \(2\).
取 \(D=(0,0,0)\)，\(B=(2,0,0)\)，\(C=(1,\sqrt{3},0)\)，\(A=(x,y,z)\).
由 \(AD=2\) 和 \(AB=a\)，得 \(x^{2}+y^{2}+z^{2}=4\)，\((x-2)^{2}+y^{2}+z^{2}=a^{2}\)，
从而 \(x=\frac{8-a^{2}}{4}\). 由 \(AC=a\) 得
\[
(x-1)^{2}+(y-\sqrt{3})^{2}+z^{2}=a^{2}
\]
结合 \(x^{2}+y^{2}+z^{2}=4\)，可得 \(y=\frac{x}{\sqrt{3}}\). 因此
\[
z^{2}=4-x^{2}-y^{2}=4-\frac{4x^{2}}{3}>0
\]
等价于 \(|x|<\sqrt{3}\)，即
\[
\sqrt{6}-\sqrt{2}<a<\sqrt{6}+\sqrt{2}
\]

再考虑两根长为 \(a\) 的棱为异棱. 设 \(AB=CD=a\)，其余四条棱长均为 \(2\).
取 \(A=\left(-\frac{a}{2},0,0\right)\)，\(B=\left(\frac{a}{2},0,0\right)\).
由 \(AC=BC=AD=BD=2\)，点 \(C\)，\(D\) 均在平面 \(x=0\) 内以原点为圆心、半径
\[
r=\sqrt{4-\frac{a^{2}}{4}}
\]
的圆上. 弦 \(CD\) 的长为 \(a\)，存在非退化弦的充要条件是 \(0<a<2r\)，即
\[
0<a<2\sqrt{2}
\]
两种情形的取值范围合并为
\[
(0,2\sqrt{2})\cup(\sqrt{6}-\sqrt{2},\sqrt{6}+\sqrt{2})=(0,\sqrt{6}+\sqrt{2})
\]
故选 A.
\end{solution}

\section{填空题}

\begin{problem}
\((1 + x + x^2)\left(x - \frac{1}{x}\right)^6\) 的展开式中的常数项为\fillinblank{}.
\end{problem}

\begin{answer}
\(-5\)
\end{answer}

\begin{solution}
由二项式定理，
\[
\left(x-\frac{1}{x}\right)^{6}=\sum_{k=0}^{6}\mathrm C_{6}^{k}(-1)^{k}x^{6-2k}
\]
乘以 \(1+x+x^{2}\) 后，常数项分别来自：\(1\) 与 \(x^{6-2k}\) 中 \(k=3\) 的项，系数为 \((-1)^{3}\mathrm C_{6}^{3}=-20\)；\(x\) 与 \(x^{6-2k}\) 无法组成常数项；\(x^{2}\) 与 \(x^{6-2k}\) 中 \(k=4\) 的项，系数为 \((-1)^{4}\mathrm C_{6}^{4}=15\). 所以常数项为 \(-20+15=-5\).
\end{solution}

\begin{problem}
已知 \(-1 < x + y < 4\) 且 \(2 < x - y < 3\)，则 \(z = 2x - 3y\) 的取值范围是\fillinblank{}. （答案用区间表示）
\end{problem}

\begin{answer}
\((3,8)\)
\end{answer}

\begin{solution}
令 \(u=x+y\)，\(v=x-y\)，则
\(x=\frac{u+v}{2}\)，\(y=\frac{u-v}{2}\).
从而
\[
z=2x-3y=\frac{5v-u}{2}
\]
因为 \(-1<u<4\)，\(2<v<3\)，所以 \(z\) 的下确界和上确界分别由 \(u\) 趋近 \(4\)、\(v\) 趋近 \(2\)，以及 \(u\) 趋近 \(-1\)、\(v\) 趋近 \(3\) 得到：
\[
\frac{5\cdot2-4}{2}=3<z<\frac{5\cdot3-(-1)}{2}=8
\]
端点均因原不等式为严格不等式而不能取到，且线性函数在这个开矩形上的值域为开区间，因此 \(z\in(3,8)\).
\end{solution}

\begin{problem}
如图，网格纸的小正方形的边长是 \(1\)，在其上用粗线画出了某多面体的三视图，则这个多面体最长的一条棱的长为\fillinblank{}.
\bitmapfigure[max width=0.36\linewidth]{img/2010/liaoning_science/q15_fig1.png}
\end{problem}

\begin{answer}
\(2\sqrt{3}\)
\end{answer}

\begin{solution}
由三视图中三条粗线的对应关系可知，同一条最长棱在三个相互垂直的投影面上的投影，均为边长 \(2\) 的正方形对角线. 设该棱在三个相互垂直的坐标方向上的分量分别为 \(u\)，\(v\)，\(w\)，则
\[
u^{2}+v^{2}=2^{2}+2^{2}=8,\qquad u^{2}+w^{2}=8,\qquad v^{2}+w^{2}=8
\]
相减可得 \(u^{2}=v^{2}=w^{2}=4\). 由空间两点间距离公式，该棱长为
\[
\sqrt{2^{2}+2^{2}+2^{2}}=2\sqrt{3}
\]
故这个多面体最长的一条棱的长为 \(2\sqrt{3}\).
\end{solution}

\begin{problem}
已知数列 \(\{a_{n}\}\) 满足 \(a_{1} = 33\)，\(a_{n+1} - a_{n} = 2n\)，则 \(\frac{a_n}{n}\) 的最小值为\fillinblank{}.
\end{problem}

\begin{answer}
\(\frac{21}{2}\)
\end{answer}

\begin{solution}
由递推关系，
\[
a_n=a_1+\sum_{k=1}^{n-1}2k=33+n(n-1)=n^{2}-n+33
\]
所以
\(\frac{a_n}{n}=n-1+\frac{33}{n}\)，且 \(n\in\mathbb Z_{+}\).
与 \(\frac{21}{2}\) 作差，得
\[
n-1+\frac{33}{n}-\frac{21}{2}=\frac{(2n-11)(n-6)}{2n}
\]
当 \(n\leqslant5\) 时，两个因式均为负，差值为正；当 \(n=6\) 时差值为 \(0\)；当 \(n\geqslant7\) 时，两个因式均为正，差值也为正. 因此最小值在 \(n=6\) 时取得，为 \(\frac{21}{2}\).
\end{solution}

\section{解答题}

\begin{problem}
在 \(\triangle ABC\) 中，\(a\)，\(b\)，\(c\) 分别为内角 \(A\)，\(B\)，\(C\) 的对边，且 \(2a \sin A = (2b + c) \sin B + (2c + b) \sin C\).

\begin{enumerate}
\item 求 \(A\) 的大小；
\item 求 \(\sin B + \sin C\) 的最大值.
\end{enumerate}
\end{problem}

\begin{solution}
\begin{enumerate}
\item 由正弦定理，设 \(a=2R\sin A\)，\(b=2R\sin B\)，\(c=2R\sin C\)，其中 \(R\) 为三角形 \(ABC\) 的外接圆半径. 将已知等式两边同时除以 \(2R\)，得
\[
2\sin ^2A=2\sin ^2B+2\sin ^2C+2\sin B\sin C
\]
即 \(\sin ^2A=\sin ^2B+\sin ^2C+\sin B\sin C\). 另一方面，由余弦定理及正弦定理，\(\sin ^2A=\sin ^2B+\sin ^2C-2\sin B\sin C\cos A\). 因为 \(B\)，\(C\) 是三角形内角，所以 \(\sin B\sin C>0\)，从而 \(-2\cos A=1\)，即 \(\cos A=-\frac{1}{2}\). 又 \(0<A<\pi\)，所以 \(A=120^{\circ}\).
\item 由 \(A+B+C=180^{\circ}\) 及第（1）问得 \(B+C=60^{\circ}\)，因此
\[
\sin B+\sin C=2\sin\frac{B+C}{2}\cos\frac{B-C}{2}=\cos\frac{B-C}{2}\leqslant1
\]
当且仅当 \(B=C=30^{\circ}\) 时等号成立，所以 \(\sin B+\sin C\) 的最大值为 \(1\).
\end{enumerate}
\end{solution}

\begin{problem}
为了比较注射 \(A\)，\(B\) 两种药物后产生的皮肤疱疹的面积，选 \(200\) 只家兔做试验，将这 \(200\) 只家兔随机地分成两组，每组 \(100\) 只，其中一组注射药物 \(A\)，另一组注射药物 \(B\). 下表 \(1\) 和表 \(2\) 分别是注射药物 \(A\) 和 \(B\) 后的试验结果. （疱疹面积单位：\(\mathrm{mm}^2\)）

\begin{center}
表 1：注射药物 \(A\) 后皮肤疱疹面积的频数分布表

\begin{tabular}{|c|c|c|c|c|}
\hline
疱疹面积 & \([60,65)\) & \([65,70)\) & \([70,75)\) & \([75,80)\) \\
\hline
频数 & 30 & 40 & 20 & 10 \\
\hline
\end{tabular}
\end{center}

\begin{center}
表 2：注射药物 \(B\) 后皮肤疱疹面积的频数分布表

\begin{tabular}{|c|c|c|c|c|c|}
\hline
疱疹面积 & \([60,65)\) & \([65,70)\) & \([70,75)\) & \([75,80)\) & \([80,85)\) \\
\hline
频数 & 10 & 25 & 20 & 30 & 15 \\
\hline
\end{tabular}
\end{center}

\begin{enumerate}
\item 甲、乙是 \(200\) 只家兔中的 \(2\) 只，求甲、乙分在不同组的概率；
\item 完成下面频率分布直方图，并比较注射两种药物后疱疹面积的中位数大小；
\item 完成下面 \(2\times2\) 列联表，并回答能否有 \(99.9\%\) 的把握认为“注射药物 \(A\) 后的疱疹面积与注射药物 \(B\) 后的疱疹面积有差异”.
\end{enumerate}

\bitmapfigure[max width=0.92\linewidth]{img/2010/liaoning_science/q18_fig1.png}

\begin{center}
\begin{tabular}{|c|c|c|c|}
\hline
 & 疱疹面积小于 \(70\,\mathrm{mm}^2\) & 疱疹面积不小于 \(70\,\mathrm{mm}^2\) & 合计 \\
\hline
注射药物 \(A\) & \(a=\)\fillinblank{} & \(b=\)\fillinblank{} &\fillinblank{} \\
\hline
注射药物 \(B\) & \(c=\)\fillinblank{} & \(d=\)\fillinblank{} &\fillinblank{} \\
\hline
合计 &\fillinblank{} &\fillinblank{} & \(n=\)\fillinblank{} \\
\hline
\end{tabular}
\end{center}

附：\(K^2=\dfrac{n(ad-bc)^2}{(a+b)(c+d)(a+c)(b+d)}\)

\begin{center}
\begin{tabular}{|c|c|c|c|c|c|}
\hline
\(P(K^2\geqslant k)\) & 0.100 & 0.050 & 0.025 & 0.010 & 0.001 \\
\hline
\(k\) & 2.706 & 3.841 & 5.024 & 6.635 & 10.828 \\
\hline
\end{tabular}
\end{center}

\end{problem}

\begin{solution}
\begin{enumerate}
\item 将甲所在的组分为第 1 组或第 2 组分别计算，得甲、乙分在不同组的概率为
\[
2\times\frac{100}{200}\times\frac{100}{199}=\frac{100}{199}
\]
\item 频率分布直方图中，矩形的高为频率除以组距. 因此，注射药物 \(A\) 后，在 \([60,65)\)，\([65,70)\)，\([70,75)\)，\([75,80)\) 上的柱高依次为 \(0.06\)，\(0.08\)，\(0.04\)，\(0.02\)；注射药物 \(B\) 后，在 \([60,65)\)，\([65,70)\)，\([70,75)\)，\([75,80)\)，\([80,85)\) 上的柱高依次为 \(0.02\)，\(0.05\)，\(0.04\)，\(0.06\)，\(0.03\). 药物 \(A\) 的累计频数达到 \(50\) 时落在区间 \([65,70)\) 内，而药物 \(B\) 的累计频数达到 \(50\) 时落在区间 \([70,75)\) 内，所以注射药物 \(A\) 后疱疹面积的中位数小于注射药物 \(B\) 后的中位数.
\item 对于药物 \(A\)，疱疹面积小于 \(70\,\mathrm{mm}^2\) 的有 \(30+40=70\) 只，不小于 \(70\,\mathrm{mm}^2\) 的有 \(20+10=30\) 只；对于药物 \(B\)，相应的两类频数分别为 \(10+25=35\) 和 \(20+30+15=65\). 因此列联表中 \(a=70\)，\(b=30\)，\(c=35\)，\(d=65\)，\(n=200\)，并且
\[
K^2=\frac{200(70\times65-30\times35)^2}{100\times100\times105\times95}\approx24.56>10.828
\]
由于 \(10.828\) 是显著性水平为 \(0.001\) 时的临界值，所以有 \(99.9\%\) 的把握认为注射药物 \(A\) 后与注射药物 \(B\) 后的疱疹面积有差异.
\end{enumerate}
\end{solution}

\begin{problem}
已知三棱锥 \(P - ABC\) 中，\(PA \perp\) 平面 \(ABC\)，\(AB \perp AC\)，\(PA = AC = \frac{1}{2} AB\)，\(N\) 为 \(AB\) 上一点，\(AB = 4AN\)，\(M\)，\(S\) 分别为 \(PB\)，\(BC\) 的中点.

\begin{enumerate}
\item 证明：\(CM \perp SN\)；
\item 求 \(SN\) 与平面 \(CMN\) 所成角的大小.
\end{enumerate}

\bitmapfigure[max width=0.36\linewidth]{img/2010/liaoning_science/q19_fig1.png}
\end{problem}

\begin{solution}
\begin{enumerate}
\item 取 \(PA=AC=1\)，则 \(AB=2\). 以 \(A\) 为原点，分别以射线 \(AB\)，\(AC\)，\(AP\) 为 \(x\)，\(y\)，\(z\) 轴的正半轴建立空间直角坐标系，则 \(A=(0,0,0)\)，\(B=(2,0,0)\)，\(C=(0,1,0)\)，\(P=(0,0,1)\)，\(N=\left(\frac{1}{2},0,0\right)\)，\(M=\left(1,0,\frac{1}{2}\right)\)，\(S=\left(1,\frac{1}{2},0\right)\). 因此 \(\overrightarrow{CM}=\left(1,-1,\frac{1}{2}\right)\)，\(\overrightarrow{SN}=\left(-\frac{1}{2},-\frac{1}{2},0\right)\)，从而
\[
\overrightarrow{CM}\cdot\overrightarrow{SN}=1\times\left(-\frac{1}{2}\right)+(-1)\times\left(-\frac{1}{2}\right)+\frac{1}{2}\times0=0
\]
所以 \(CM\perp SN\).
\item 由 \(\overrightarrow{CN}=\left(\frac{1}{2},-1,0\right)\)，得平面 \(CMN\) 的一个法向量可以取为
\[
\bs n=4\left(\overrightarrow{CM}\times\overrightarrow{CN}\right)=(2,1,-2)
\]
令 \(\bs d=\overrightarrow{SN}=\left(-\frac{1}{2},-\frac{1}{2},0\right)\)，设 \(SN\) 与平面 \(CMN\) 所成的角为 \(\theta\)，则
\[
\sin\theta=\frac{|\bs d\cdot\bs n|}{|\bs d|\,|\bs n|}=\frac{\frac{3}{2}}{\frac{1}{\sqrt{2}}\times3}=\frac{\sqrt{2}}{2}
\]
因为线面角取锐角或直角，所以 \(\theta=45^{\circ}\).
\end{enumerate}
\end{solution}

\begin{problem}
设椭圆 \(C: \frac{x^{2}}{a^{2}} + \frac{y^{2}}{b^{2}} = 1\) （\(a > b > 0\)）的左焦点为 \(F\)，过点 \(F\) 的直线与椭圆 \(C\) 相交于 \(A\)，\(B\) 两点，直线 \(l\) 的倾斜角为 \(60^\circ\)，\(\overrightarrow{AF} = 2\overrightarrow{FB}\).

\begin{enumerate}
\item 求椭圆 \(C\) 的离心率；
\item 如果 \(\left|AB\right| = \frac{15}{4}\)，求椭圆 \(C\) 的方程.
\end{enumerate}
\end{problem}

\begin{solution}
\begin{enumerate}
\item 设椭圆的半焦距为 \(c\)，则 \(c^2=a^2-b^2\)，左焦点为 \(F=(-c,0)\). 取倾斜角为 \(60^{\circ}\) 的直线方向单位向量 \(\left(\frac{1}{2},\frac{\sqrt{3}}{2}\right)\)，则直线 \(l\) 上的点可表示为
\[
(x,y)=\left(-c,0\right)+t\left(\frac{1}{2},\frac{\sqrt{3}}{2}\right)
\]
即 \(x=-c+\frac{t}{2}\)，\(y=\frac{\sqrt{3}t}{2}\). 设点 \(A\)，\(B\) 对应的参数分别为 \(t_A\)，\(t_B\)，由 \(\overrightarrow{AF}=2\overrightarrow{FB}\) 得 \(t_A=-2t_B\). 将直线参数方程代入椭圆方程，得
\[
\frac{\left(-c+\frac{t}{2}\right)^2}{a^2}+\frac{3t^2}{4b^2}=1
\]
整理得
\[
\left(b^2+3a^2\right)t^2-4b^2ct-4b^4=0
\]
记 \(D=b^2+3a^2\)，并令 \(q=-t_B>0\)，则 \(t_A=2q\). 由韦达定理，\(t_A+t_B=q=\frac{4b^2c}{D}\)，且 \(t_At_B=-2q^2=-\frac{4b^4}{D}\)，所以 \(q^2=\frac{2b^4}{D}\). 将 \(q=\frac{4b^2c}{D}\) 代入，得 \(D=8c^2\)，即 \(b^2+3a^2=8(a^2-b^2)\)，从而 \(9b^2=5a^2\). 因此
\[
e=\frac{c}{a}=\sqrt{1-\frac{b^2}{a^2}}=\frac{2}{3}
\]
\item 因为所取方向向量为单位向量，所以 \(|AB|=|t_A-t_B|=3q\). 由 \(D=8c^2\) 及 \(q=\frac{4b^2c}{D}\)，得 \(q=\frac{b^2}{2c}\)，从而
\[
\frac{15}{4}=|AB|=\frac{3b^2}{2c}
\]
即 \(b^2/c=5/2\). 由上问 \(b^2=5a^2/9\)，\(c=2a/3\)，所以 \(b^2/c=5a/6=5/2\)，解得 \(a=3\)，\(b^2=5\). 故椭圆 \(C\) 的方程为 \(\dfrac{x^2}{9}+\dfrac{y^2}{5}=1\).
\end{enumerate}
\end{solution}

\begin{problem}
已知函数 \( f(x) = (a + 1)\ln x + ax^2 + 1 \).

\begin{enumerate}
\item 讨论函数 \( f(x) \) 的单调性；
\item 设 \(a < -1\)，如果对任意 \(x_{1}\)，\(x_{2} \in (0, +\infty)\)，\(|f(x_{1}) - f(x_{2})| \geqslant 4|x_{1} - x_{2}|\)，求 \(a\) 的取值范围.
\end{enumerate}
\end{problem}

\begin{solution}
\begin{enumerate}
\item 函数的定义域为 \((0,+\infty)\)，且
\[
f'(x)=\frac{a+1}{x}+2ax=\frac{a+1+2ax^2}{x}
\]
因为 \(x>0\)，只需讨论 \(a+1+2ax^2\) 的符号. 当 \(a\geqslant0\) 时，\(a+1+2ax^2>0\)，所以 \(f(x)\) 在 \((0,+\infty)\) 上单调增加；当 \(a\leqslant-1\) 时，\(a+1+2ax^2<0\)，所以 \(f(x)\) 在 \((0,+\infty)\) 上单调减少；当 \(-1<a<0\) 时，令 \(f'(x)=0\)，得 \(x_0=\sqrt{-\dfrac{a+1}{2a}}\)，故 \(f(x)\) 在 \(\left(0,\sqrt{-\dfrac{a+1}{2a}}\right)\) 上单调增加，在 \(\left(\sqrt{-\dfrac{a+1}{2a}},+\infty\right)\) 上单调减少.
\item 题设给出 \(a<-1\)，此时 \(f'(x)<0\). 对任意不同的 \(x_1\)，\(x_2\)，由拉格朗日中值定理，存在 \(\xi\) 严格介于二者之间，使得
\[
\frac{|f(x_1)-f(x_2)|}{|x_1-x_2|}=|f'(\xi)|
\]
所以若 \(|f'(x)|\geqslant4\) 对任意 \(x>0\) 成立，则由中值定理可知题设不等式成立；反过来，若题设不等式对任意 \(x_1\)，\(x_2\in(0,+\infty)\) 成立，令 \(x_2\) 趋于 \(x_1=x\)，可得 \(|f'(x)|\geqslant4\). 因此二者等价. 又
\[
|f'(x)|=-f'(x)=\frac{-a-1}{x}-2ax
\]
由基本不等式，\(\frac{-a-1}{x}-2ax\geqslant2\sqrt{(-a-1)(-2a)}=2\sqrt{2a(a+1)}\)，且等号可在正数 \(x=\sqrt{\frac{-a-1}{-2a}}\) 时取得. 所以所需条件为 \(2\sqrt{2a(a+1)}\geqslant4\)，即 \(a(a+1)\geqslant2\). 结合 \(a<-1\)，解得 \(a\leqslant-2\).
\end{enumerate}
\end{solution}

\section{选考题：解答题}

\begin{problem}
如图，\(\triangle ABC\) 的角平分线 \(AD\) 的延长线交它的外接圆于点 \(E\).

\begin{enumerate}
\item 证明：\(\triangle ABE \sim \triangle ADC\)；
\item 若 \(\triangle ABC\) 的面积 \(S = \frac{1}{2} AD \cdot AE\)，求 \(\angle BAC\) 的大小.
\end{enumerate}

\bitmapfigure[max width=0.34\linewidth]{img/2010/liaoning_science/q22_fig1.png}
\end{problem}

\begin{solution}
\begin{enumerate}
\item 因为 \(AD\) 是 \(\angle BAC\) 的平分线，且 \(A\)，\(D\)，\(E\) 共线，所以 \(\angle BAE=\angle CAD\). 又点 \(A\)，\(B\)，\(C\)，\(E\) 共圆，故同弧 \(AB\) 所对的圆周角相等，即 \(\angle AEB=\angle ACB\). 由于 \(B\)，\(D\)，\(C\) 共线，所以 \(\angle ACB=\angle ACD\). 因此由两角对应相等，得 \(\triangle ABE\sim\triangle ADC\).
\item 由相似三角形的对应边成比例，得 \(\dfrac{AB}{AD}=\dfrac{AE}{AC}\)，即 \(AB\cdot AC=AD\cdot AE\). 另一方面，\(\triangle ABC\) 的面积为
\[
S=\frac{1}{2}AB\cdot AC\sin\angle BAC
\]
结合题设 \(S=\frac{1}{2}AD\cdot AE\) 及上式，得 \(\sin\angle BAC=1\). 因为 \(\angle BAC\) 是三角形内角，所以 \(\angle BAC=90^{\circ}\).
\end{enumerate}
\end{solution}

\begin{problem}
已知 \(P\) 为半圆 \(C:\begin{cases}
x = \cos \theta,\\
y = \sin \theta.
\end{cases}\) （\(\theta\) 为参数，\(0 \leqslant \theta \leqslant \pi\)）上的点，点 \(A\) 的坐标为 \((1,0)\)，\(O\) 为坐标原点，点 \(M\) 在射线 \(OP\) 上，线段 \(OM\) 与 \(C\) 的弧 \(\widehat{AP}\) 的长度均为 \(\frac{\pi}{3}\).

\begin{enumerate}
\item 以 \(O\) 为极点，\(x\) 轴的正半轴为极轴建立极坐标系，求点 \(M\) 的极坐标；
\item 求直线 \(AM\) 的参数方程.
\end{enumerate}
\end{problem}

\begin{solution}
\begin{enumerate}
\item 由参数方程可知半圆 \(C\) 是以 \(O\) 为圆心、半径为 \(1\) 的上半圆，点 \(A\) 对应参数 \(\theta=0\). 由于弧 \(\widehat{AP}\) 的长度等于 \(\frac{\pi}{3}\)，且半径为 \(1\)，所以弧 \(\widehat{AP}\) 所对的圆心角为 \(\frac{\pi}{3}\)，从而点 \(P\) 对应的参数为 \(\theta=\frac{\pi}{3}\). 因此射线 \(OP\) 的极角为 \(\frac{\pi}{3}\). 又因为 \(OM=\frac{\pi}{3}\)，所以点 \(M\) 的极坐标为 \(\left(\frac{\pi}{3},\frac{\pi}{3}\right)\).
\item 由点 \(M\) 的极坐标可得 \(M\) 的直角坐标为 \(\left(\frac{\pi}{3}\cos\frac{\pi}{3},\frac{\pi}{3}\sin\frac{\pi}{3}\right)=\left(\frac{\pi}{6},\frac{\sqrt{3}\pi}{6}\right)\). 因此 \(\overrightarrow{AM}=\left(\frac{\pi}{6}-1,\frac{\sqrt{3}\pi}{6}\right)\)，取其 \(6\) 倍向量 \((\pi-6,\sqrt{3}\pi)\) 作为直线 \(AM\) 的方向向量，得直线 \(AM\) 的参数方程为：
\[
\begin{cases}
x=1+(\pi-6)t,\\
y=\sqrt{3}\pi t,
\end{cases}\qquad t\in\R
\]
\end{enumerate}
\end{solution}

\begin{problem}
已知 \(a\)，\(b\)，\(c\) 均为正数，证明：\(a^2 + b^2 + c^2 + \left(\frac{1}{a} + \frac{1}{b} + \frac{1}{c}\right)^2 \geqslant 6\sqrt{3}\)，并确定 \(a\)，\(b\)，\(c\) 为何值时，等号成立.
\end{problem}

\begin{solution}
因为 \(a\)，\(b\)，\(c\) 均为正数，由均值不等式得 \(a^2+b^2+c^2\geqslant3\sqrt[3]{a^2b^2c^2}=3(abc)^{\frac{2}{3}}\)，且 \(\frac{1}{a}+\frac{1}{b}+\frac{1}{c}\geqslant\frac{3}{\sqrt[3]{abc}}\)，所以 \(\left(\frac{1}{a}+\frac{1}{b}+\frac{1}{c}\right)^2\geqslant\frac{9}{(abc)^{\frac{2}{3}}}\). 令 \(u=(abc)^{\frac{2}{3}}>0\)，则原式满足：
\[
a^2+b^2+c^2+\left(\frac{1}{a}+\frac{1}{b}+\frac{1}{c}\right)^2\geqslant3u+\frac{9}{u}\geqslant2\sqrt{27}=6\sqrt{3}
\]
当且仅当 \(a=b=c\) 且 \(3u=\frac{9}{u}\) 时等号成立. 由 \(u=(abc)^{\frac{2}{3}}\) 和 \(a=b=c\)，设 \(a=b=c=t>0\)，则 \(u=t^2\)，从而 \(3t^2=\frac{9}{t^2}\)，即 \(t^4=3\)，所以 \(t=\sqrt[4]{3}\). 故当且仅当 \(a=b=c=\sqrt[4]{3}\) 时，等号成立.
\end{solution}
